\chapter{射影簇之间的映射}\label{chap:maps-of-projective-varieties}

\noindent\textbf{第 \thechapter{} 章学习目标}
\begin{itemize}
  \item 定义并研究射影簇之间的正则映射.
  \item 举出正则映射的各种例子, 包括线性映射, 多项式映射, 同构, 射影等价和 Veronese 嵌入.
  \item 定义 Segre 映射, 并用它们定义射影簇的积.
  \item 抽象地描述 Grassmann 簇, 并通过 Plücker 映射将它们实现为射影簇.
\end{itemize}

在第~\ref{chap:polynomial-maps}~章的引言中, 我们鼓励读者每逢引入一类新的数学对象, 就问一个关键问题: 这些对象之间哪些映射保持所关心的结构? 对于仿射簇, 最终得到适当的保结构映射概念是多项式映射. 本章的第一个目标是在射影情形中定义相应的概念.

仿射簇和射影簇之间若干关键的差别使这一目标变得复杂. 首先, 由于 $\mathbb{P}^{m}$ 定义中的等价关系, 一个多项式 $f\in K[x_{0},\ldots,x_{m}]$ 并不给出射影簇 $X\subseteq\mathbb{P}^{m}$ 上的良定义函数, 原因仅仅是多项式的值会随 $\mathbb{P}^{m}$ 中齐次坐标的缩放而改变. 不过, 若 $f$ 恰是 $d$ 次齐次多项式, 我们已经看到, 缩放齐次坐标对 $f$ 的值有一个相当简单的影响:
\[
  f(\lambda a_{0},\ldots,\lambda a_{m})=\lambda^{d}f(a_{0},\ldots,a_{m})
  \qquad\text{对所有 }\lambda\in K\setminus\{0\}.
\]
因此, 给定同为一个次数的齐次多项式 $f_{0},\ldots,f_{n}\in K[x_{0},\ldots,x_{m}]$, 其中至少有一个不在 $a$ 消没, 就得到良定义的值 $[f_{0}(a):\cdots:f_{n}(a)]\in\mathbb{P}^{n}$, 它不依赖于 $a$ 的齐次坐标的选取. 受这个观察启发, 我们以这样的多项式组为基础研究射影簇之间的映射.

一般来说, 若 $X\subseteq\mathbb{P}^{m}$ 和 $Y\subseteq\mathbb{P}^{n}$ 为射影簇, 则称函数 $F:X\to Y$ 为\emph{正则映射}, 如果它至少局部可以由一组同次数的齐次多项式 $f_{0},\ldots,f_{n}$ 实现. 正则映射是射影簇之间的保结构映射, 本章前半部分的首要目标是熟悉它们.

本章后半部分描述两种重要构造: 射影簇的积和 Grassmann 簇. 起初, 这两种构造都并不显然给出射影簇; 但对于每一种, 我们都会展示能赋予它射影簇结构的自然映射 (分别是 Segre 映射和 Plücker 映射).

\par\smallskip
{\footnotesize
\noindent\textcopyright{} The Author(s) 2025\par
\noindent E. Clader, D. Ross, \emph{Beginning in Algebraic Geometry}, Undergraduate Texts in Mathematics,
\href{https://doi.org/10.1007/978-3-031-88819-9_11}{\nolinkurl{https://doi.org/10.1007/978-3-031-88819-9_11}}.\par}
